% !TeX root = ../main.tex
\section{Kế hoạch gọi vốn}

% !TeX root = ../main.tex
\subsection{Nhu cầu gọi vốn}

\textbf{1. Xác định nhu cầu vốn đầu tư:}
Dựa trên mô hình dự phóng tài chính và phân tích dòng tiền ở Phần 4, khoản thiếu hụt dòng tiền tích lũy lớn nhất (Maximum Cash Drain) của Mutux rơi vào \textbf{Quý 3 năm đầu tiên} với mức âm tích lũy là \textbf{-231 triệu VNĐ}. Mặc dù sang Quý 4 dòng tiền thuần từ hoạt động kinh doanh bắt đầu đạt dương (+\textbf{91 triệu VNĐ}), tổng khoản thiếu hụt tài chính tích lũy toàn năm 1 vẫn ở mức khoảng \textbf{140 - 156 triệu VNĐ}.

Nhằm đảm bảo an toàn tài chính, chuẩn bị quỹ dự phòng cho rủi ro vận hành (rủi ro thanh khoản cọc, rủi ro nợ xấu tài sản) và duy trì thời gian hoạt động (Runway) tối thiểu 18 tháng, Mutux xác định nhu cầu gọi vốn cho vòng Tiền hạt giống (Pre-Seed Round) là \textbf{12.000 USD} (tương đương khoảng \textbf{300.000.000 VNĐ}).

\textbf{2. Mục tiêu chiến lược của vòng gọi vốn Pre-Seed:}
\begin{itemize}
    \item \textbf{Đảm bảo Runway hoạt động:} Cung cấp nguồn vốn an toàn kéo dài ít nhất 18--24 tháng, giúp startup vượt qua giai đoạn khởi tạo (Cold-start) và duy trì bộ máy vận hành tinh gọn.
    \item \textbf{Đạt điểm hòa vốn (Break-even Point):} Tài trợ nguồn lực để thúc đẩy tốc độ tăng trưởng giao dịch, đưa nền tảng đạt điểm hòa vốn vào \textbf{năm thứ 2 (Quý 3)} khi doanh thu định kỳ đạt mức 80--100 triệu VNĐ/tháng.
    \item \textbf{Tạo đà bứt phá quy mô:} Đạt mục tiêu 10.000 người dùng đăng ký, 7.050 đơn thuê thành công và 1,11 tỷ VNĐ tổng giá trị giao dịch (GMV) trong 12 tháng đầu, làm tiền đề mở rộng các dịch vụ B2B và thương mại hóa chuyên sâu.
\end{itemize}


% !TeX root = ../main.tex
\subsection{Mục đích sử dụng vốn}

Nguồn vốn gọi được từ vòng Pre-Seed (300.000.000 VNĐ) sẽ được tập trung đầu tư vào 3 trụ cột chiến lược chính nhằm giải quyết các rủi ro trọng yếu của mô hình Marketplace cho thuê:

\begin{enumerate}
    \item \textbf{Vượt qua rào cản Cold-start \& Chi phí thu hút khách hàng (CAC):}
    \begin{itemize}
        \item \textit{Marketing kĩ thuật số (Digital Performance Marketing):} Triển khai các chiến dịch quảng cáo nhắm mục tiêu trên TikTok, Facebook và Google Ads đến tập khách hàng sinh viên, game thủ và người yêu công nghệ.
        \item \textit{Chương trình khuyến mãi \& Trợ giá đơn đầu:} Đội ngũ chấp nhận bù lỗ chi phí đơn thuê đầu tiên thông qua các mã giảm giá/voucher nhằm giảm tối đa rào cản thử nghiệm dịch vụ của người dùng mới.
        \item \textit{Kênh truyền thông trường đại học \& E-sports Community:} Đồng hành và tài trợ các giải đấu E-sports sinh viên, hợp tác với các Câu lạc bộ Công nghệ/Game tại các trường đại học lớn trên địa bàn TP.HCM.
    \end{itemize}

    \item \textbf{Xây dựng \& Hoàn thiện hạ tầng công nghệ lõi (Tech \& Infrastructure) - hoàn thiện từ MVP thành sản phẩm thật:}
    \begin{itemize}
        \item \textit{Cổng thanh toán Escrow tự động \& PayOS:} Hoàn thiện cơ chế giữ và giải ngân tiền thuê tự động, tích hợp cổng đối soát PayOS giúp xử lý giao dịch tức thì và chính xác.
        \item \textit{Xác thực danh tính eKYC:} Tích hợp giải pháp eKYC tự động nhằm xác minh định danh người thuê, giảm thiểu tối đa rủi ro lừa đảo, tráo đổi hoặc mất mát thiết bị gaming gear.
        \item \textit{Tích hợp dịch vụ đối tác tài chính (BNPL):} Kết nối API mượt mà với các nền tảng BNPL (MoMo, Fundiin) để vận hành Dịch vụ Hỗ trợ tiền cọc (Deposit Advance Service), giảm bớt rào cản đặt cọc cho người thuê.
        \item \textit{Duy trì Cloud Server \& API Services:} Chi trả phí hạ tầng máy chủ đám mây, dịch vụ tin nhắn SMS OTP và tên miền nền tảng.
    \end{itemize}

    \item \textbf{Duy trì bộ máy vận hành tinh gọn (Lean Operations) \& Quỹ dự phòng rủi ro:}
    \begin{itemize}
        \item \textit{Chi phí nhân sự nòng cốt:} Trang trải phụ cấp và lương cơ bản cho Đội ngũ sáng lập, lập trình viên lõi và nhân viên Chăm sóc khách hàng (CSKH).
        \item \textit{Xử lý tranh chấp \& Kiểm định thiết bị:} Duy trì quy trình tiếp nhận, kiểm tra chất lượng thiết bị và giải quyết khiếu nại giữa Bên cho thuê (Lender) và Bên thuê (Renter).
        \item \textit{Tư vấn Pháp lý \& Quỹ dự phòng:} Chi phí hoàn thiện Điều khoản dịch vụ (ToS), hợp đồng điện tử và trích lập quỹ dự phòng xử lý nợ xấu hoặc hỏng hóc tài sản phát sinh.
    \end{itemize}
\end{enumerate}


% !TeX root = ../main.tex
\subsection{Kế hoạch sử dụng vốn theo từng hạng mục và lộ trình giải ngân}

\textbf{1. Bảng phân bổ nguồn vốn 300.000.000 VNĐ (12.000 USD):}
Nguồn vốn gọi được sẽ được cam kết phân bổ tuân thủ nghiêm ngặt theo tỷ lệ ngân sách chi tiết như sau:

\begin{table}[htbp]
    \centering
    \caption{Phân bổ chi tiết nguồn vốn gọi được (Pre-Seed Round)}
    \resizebox{\textwidth}{!}{%
    \begin{tabular}{@{}lp{0.45\textwidth}rr@{}}
        \toprule
        \textbf{Hạng mục đầu tư} & \textbf{Nội dung chi tiết} & \textbf{Số tiền (VNĐ)} & \textbf{Tỷ lệ (\%)} \\
        \midrule
        \textbf{Marketing \& CAC} & Digital Ads, voucher trợ giá đơn đầu, event sinh viên, CLB E-sports. & 90.000.000 & 30,0\% \\
        \textbf{Phát triển Sản phẩm \& Hạ tầng} & Tích hợp Escrow, eKYC, PayOS, kết nối BNPL, Cloud Server \& API. & 80.000.000 & 26,7\% \\
        \textbf{Vận hành \& Nhân sự} & Phụ cấp Team sáng lập, đội ngũ CSKH, kiểm định \& xử lý tranh chấp. & 90.000.000 & 30,0\% \\
        \textbf{Tư vấn Pháp lý} & Hợp đồng điện tử, điều khoản dịch vụ, bản quyền \& tư vấn pháp lý. & 20.000.000 & 6,7\% \\
        \textbf{Quỹ Dự phòng} & Dự phòng nợ xấu, tổn thất tài sản thiết bị \& rủi ro thanh khoản cọc. & 20.000.000 & 6,7\% \\
        \midrule
        \textbf{TỔNG CỘNG} & & \textbf{300.000.000} & \textbf{100,0\%} \\
        \bottomrule
    \end{tabular}%
    }
\end{table}

\textbf{2. Lộ trình giải ngân theo 4 Quý năm đầu:}
Để quản trị rủi ro dòng tiền và tối ưu hiệu quả sử dụng vốn, nguồn vốn 300 triệu VNĐ sẽ được giải ngân theo 4 giai đoạn tương ứng với các cột mốc phát triển của dự án:

\begin{itemize}
    \item \textbf{Quý 1 (Giai đoạn Triển khai MVP \& Khởi động - Giải ngân 40\% $\sim$ 120 triệu VNĐ):}
    Tập trung 60 triệu cho phát triển công nghệ lõi (Escrow, eKYC, PayOS), 30 triệu hoàn thiện pháp lý/hợp đồng điện tử và 30 triệu khởi động chiến dịch Marketing Cold-start.
    
    \item \textbf{Quý 2 (Giai đoạn Thử nghiệm \& Tăng trưởng - Giải ngân 30\% $\sim$ 90 triệu VNĐ):}
    Giải ngân 40 triệu cho Marketing kĩ thuật số (TikTok/Facebook Ads) và trợ giá đơn thuê; 30 triệu duy trì vận hành CSKH; 20 triệu nâng cấp tính năng nền tảng.
    
    \item \textbf{Quý 3 (Giai đoạn Bứt phá Quy mô - Giải ngân 20\% $\sim$ 60 triệu VNĐ):}
    Phân bổ 30 triệu cho các chương trình Marketing sự kiện sinh viên/Gaming Cafe; 20 triệu cho nhân sự vận hành; 10 triệu duy trì hạ tầng máy chủ.
    
    \item \textbf{Quý 4 (Giai đoạn Chuẩn bị Hòa vốn \& Mở rộng B2B - Giải ngân 10\% $\sim$ 30 triệu VNĐ):}
    Dành 15 triệu hoàn thiện tính năng hiển thị ưu tiên/đẩy tin cho đối tác B2B; 15 triệu trích lập bổ sung Quỹ dự phòng tài sản trước khi bước sang Năm 2.
\end{itemize}

\textbf{3. Chỉ số kiểm soát hiệu quả sử dụng vốn (KPIs):}
Hiệu quả giải ngân vốn sẽ được đo lường chặt chẽ hàng tháng thông qua các chỉ số hiệu suất:
\begin{itemize}
    \item \textbf{Chi phí thu hút người dùng thanh toán (CAC):} Giữ mức CAC hiệu dụng dưới 60.000 VNĐ/user thuê thực tế.
    \item \textbf{Tỷ lệ nợ xấu/thất thoát tài sản:} Đảm bảo tỷ lệ tổn thất thiết bị dưới \textbf{1\%} tổng số đơn thuê nhờ cơ chế eKYC + Escrow.
    \item \textbf{Tốc độ quay vòng vốn cọc (Deposit Velocity):} Tỷ lệ chuyển đổi người dùng đăng ký mở ví cọc đạt tối thiểu 30\% tổng số người thuê thực tế.
\end{itemize}


